\section{Mechanization and the limit-periodic corollary}\label{sec:mechanization}

\subsection{What is checked by machine}\label{sec:lean}

% ledger: M1, M2, T12, D4, N8; source lean/R44/AXIOMS.md, HYPOTHESES.md, THEOREM.md at the pinned commit; profile §7
The finite parts of the proof are theorems of a Lean~4 development written for this paper, with
Mathlib as its library, and the logical assembly of the whole argument is a Lean theorem as
well. \Cref{tab:finite} lists the $23$ finite theorems with their statements and the method by
which each is decided: by the kernel's evaluator (\lean{decide}) where that is feasible, and
otherwise by \lean{native\_decide}, which compiles the decision procedure and trusts its result
through one named axiom per theorem. Those axioms are the only non-standard ones in the
development: \cref{tab:axioms} gives, for every finite theorem and for the main theorems, the
output of Lean's \lean{\#print axioms}, which for the main theorem consists of the three standard
axioms of classical Mathlib and the $21$ named hooks of the finite checks it uses. No theorem in
the development depends on \lean{sorryAx}, and no admission remains.

% ledger: T12, M2b (closed), D4; the honest sentence; the Hypotheses record (empty)
The assembly is the theorem \lean{r44\_einstein}, which takes no hypothesis: it states that a
tiling exists, that every tiling has $\Per(T)=\{0\}$, and that $|\Sym(T)|\le24$; the theorems
\lean{tiling\_homochiral} and \lean{carrier\_hierarchy}, with
\lean{geometric\_hierarchy\_unique}, state clauses (3) and (4) of \cref{thm:full}. The
development was written with a record \lean{Hypotheses} of named written lemmas as its
verification boundary, and the written lemmas of this paper are the fields that record had: the
tube construction and carrier recovery in \cref{app:tubes,app:asymmetry}, realization in \cref{sec:existence},
and the companion and registration arguments in \cref{sec:companions,sec:registration}. Every one
of them is now a Lean theorem of the pinned development (\cref{tab:hypotheses} lists the $19$
theorems, with their tiers by their own axiom lines), the record is empty, and the conditional
form survives under the name \lean{r44\_einstein\_of\_hypotheses} over the empty record
(\cref{app:lean}). Everything else, including the period-halving tower, the symmetry bound, local
finiteness, the unique parent, the coarsening and the hierarchy with its uniqueness among
arbitrary geometrically nested registered families (\cref{thm:hierarchy}), is proved inside
Lean. The honest summary is therefore this: the theorem is kernel-checked modulo the named
compiler hooks; the written proofs of \cref{app:solidchecks,sec:existence} and
\cref{sec:companions,sec:registration} remain as exposition, independently reviewed, and four of
the formal proofs argue differently from the written ones, as noted where those lemmas appear
(carrier recovery, containment rigidity, the covering by a component's solids, and small-collar
realization). The same pipeline was exercised on the hat as a calibration of the tooling
(\cref{app:calibration}). The hat and Spectre papers rely on no proof assistant; the nearest comparator is Myers's
in-progress formalization of those papers \cite{Mye24lean}, and we claim no coverage by Mathlib
itself \cite{Lean4}.

% ledger: M3, M4, M5, M8; source APPENDIX_DATA.md; verify/replay.py; controls.sh
Independently of Lean, every enumeration is replayed from the Python standard library by one
command in about a minute (\cref{app:data}), and the companion census was checked by two
independent Python implementations; the second extracts the features from the actual
triangulated mesh rather than from their annotation table. No floating-point value enters any
checked statement: certificates are integers and exact rationals, the checkers use exact
arithmetic, and the Lean theorems decide integer literals. Negative controls accompany the
positive build. The diagonal audit distinguishes assertions about distinct companion tiles from
relations with legitimate reflexive cases; the controls exercise the repaired diagonal
application, corrupted finite inputs, mixed-handedness claims, a second-hierarchy claim and a
claim that a raw geometric hierarchy has a noncanonical level-one pose, each
checked to fail with its expected diagnostic. Six corrupted inputs to the replay are rejected. The
development was reviewed adversarially twice before this paper was written; the
first round found the hypothesis record inconsistent (one field admitted comparing a tile with
itself), which was repaired and controlled, and the second round found no blocking or major
issue, while noting that the reviewer could not run Lean. A third round, on the closed record with
the empty hypothesis record, was requested on 2026-09-10; as of that date its verdict is
pending. We make no claim of simplicity, and we
have one route to the theorem: the continuous argument of \cref{sec:companions} is written for
checking by hand with its arithmetic in Lean, and the finite part is where the computer does
the work.

\emph{Code.} The repository containing the solid, the certificates, the replay, the Lean
development and the viewer accompanies this paper (\cref{app:data}); its mathematical baseline is commit
\texttt{d90313a717}, the intended release tag is \texttt{paper-v1}, and tag creation and licensing
remain pending owner actions.

% ledger: D4, T12, M2, M6, M7; the freeze (step 18) and the re-pin (step 19)
\emph{Version freeze.} This review version uses commit \texttt{d90313a717} (2026-09-10) as its
mathematical and Lean baseline. The $19$ theorems of \cref{tab:hypotheses}, the empty record, and
the recorded axiom dependencies of \cref{tab:axioms} refer to that baseline; the geometric
certificate data and Lean statements refer to that baseline, and the bibliography and replay
measurements have the separately identified sources recorded in the paper. The intended
paper-release tag is \texttt{paper-v1}, to be placed on the paper-release commit; tag creation
and licensing remain pending owner actions. The theorem, the written proof and the finite
certificates are unchanged since the earlier baseline \texttt{dd2735d9bd} (2026-09-08), at which
$16$ of the $19$ written lemmas were still hypotheses of the Lean assembly.

% ledger: N9; simulations/periodicity_search.py; simulations/periodicity_search_report.json (commit d90313a717); simulations/review/2026-09-09_periodicity_search_r1_eddy.md (controls run)
\emph{Evidence, not proof.} Independently of the proof, and used nowhere in it, a theorem-free
search was run over grid-registered copies of $\solid$ (all $48$ signed frames, reflections
included; copies compatible when their carrier cells are disjoint and the features on every
shared unit panel are exact opposites, in rational arithmetic;
\file{simulations/periodicity\_search.py}). A SAT solver was asked, for every full-rank sublattice
$L\subset\Z^3$ of index at most $40$, whether the torus $\R^3/L$ can be tiled by such copies;
all $43{,}981$ instances are unsatisfiable, so no tiling by registered copies has a full-rank
period lattice of index at most $40$. Heesch-style coronas were grown around one copy to radius
$5$ of $5$ tried, with $251$ copies at radius $5$. The positive controls of the same script find
the unit cube periodic at index $1$ and the featureless chair carrier periodic at index $7$,
as recorded in \file{simulations/review/2026-09-09_periodicity_search_r1_eddy.md}.
The torus and corona receipt is \file{simulations/periodicity\_search\_report.json}
(wall time $108{,}682$ s). This is
evidence in the sense of a sanity check on the finite data: it concerns registered copies and
full-rank lattices only, whereas \cref{thm:noperiod} excludes every period of every tiling.

\input{generated/finite_theorems_table}

% generated by paper/scripts/hypotheses_table.py; from lean/R44/HYPOTHESES.md 'Discharged', PAPER_UPDATES.md section 2 and build_axioms.log at d90313a717; 19 former fields; 0 remaining
{\small
\begin{longtable}{L{0.21\linewidth} L{0.09\linewidth} L{0.22\linewidth} L{0.05\linewidth} L{0.31\linewidth}}
\caption{The written lemmas of Sections 2--5 that formed the \lean{Hypotheses} record, each now a Lean theorem of the pinned development (the record is empty, Appendix C): the former field, the written lemma it quoted, the theorem, its tier by its own axiom line, and the argument of the formal proof where it differs from, or refines, the written one (FIGURES.md, Table 3).}\label{tab:hypotheses}\\
\toprule
former field & written lemma & Lean theorem & tier & how the formal proof argues \\
\midrule
\endfirsthead
\toprule
former field & written lemma & Lean theorem & tier & how the formal proof argues \\
\midrule
\endhead
\bottomrule
\endfoot
\lean{retained\_core\_overlap} & A-L5.1 & \lean{retained\_core\_overlap\_holds} & T1 & density choice of overlapping cube interiors; eighth-grid midpoint \\
\lean{carrier\_feature\_frame\_reduction} & R\S{}6, E5 & \lean{carrier\_feature\_frame\_reduction\_holds} & T1n & six carrier symmetries from diameter pairs and the missing corner; signed heights from membership thresholds along normals \\
\lean{circular\_cone\_solid\_angle} & A-L4.1 & \lean{circular\_cone\_solid\_angle\_holds} & T1 & horizontal disk slices; volume (2$\pi$/\allowbreak{}3)(1 - L/\allowbreak{}$\surd$(1+L$^2$)), L = 3/\allowbreak{}25 \\
\lean{per\_tube\_homeomorphisms} & A\S{}1 & \lean{per\_tube\_homeomorphisms\_holds} & T1n & inverse written with the continuous clamp $\psi$; seams and joint continuity in the height proved \\
\lean{feature\_tube\_maps\_glue} & A\S{}1 & \lean{feature\_tube\_maps\_glue\_holds} & T1n & gluing by finite composition of supported homeomorphisms \\
\lean{feature\_tube\_map\_carries\_carrier} & A\S{}1 & \lean{feature\_tube\_map\_carries\_carrier\_holds} & T1n & F(x) $\in$ Q $\Leftrightarrow$ x $\in$ P proved pointwise, not inferred from retained cores \\
\lean{feature\_circular\_cone\_containment} & A-L4.1 & \lean{feature\_circular\_cone\_containment\_holds} & T1n & agreement with the signed-tent hypograph on an open ambient neighbourhood of every feature-graph point (base edges included); Lipschitz cone inclusion transported by an affine isometry \\
\lean{planar\_area\_carrier\_recovery} & R\S{}6, E5 & \lean{planar\_area\_carrier\_recovery\_holds} & T1 & alternative proof of the unchanged proposition: diameter-endpoint argument (squared distance bound 1089/\allowbreak{}100 < 12; the six antipodal carrier corners recover centre and axes; the missing corner excludes reversal). Not the written planar-area calculation; say so where the paper cites R\S{}6 \\
\lean{cone\_sector\_budgets} & A-L2.2 & \lean{cone\_sector\_budgets\_holds} & T1n & Q as a finite Boolean combination of affine halfspaces; uniform conical germs; equality with the radial tangent cone; coverage includes cone-boundary directions, disjointness is of interiors \\
\lean{connected\_feature\_companion} & A-T4.2 & \lean{connected\_feature\_companion\_holds} & T1n & finite disjoint closed cover on connected sets; triple contact excluded via the discharged open-cone inclusion and volume formula (no assumption that a cone and its interior have equal volume) \\
\lean{feature\_containment\_rigidity} & A-L4.3, E1 & \lean{feature\_containment\_rigidity\_holds} & T1n & compactness argument (K compact, f an isometry, f(K) $\subseteq$ K $\Rightarrow$ f(K) = K) in place of the written graph-length argument; uses the distinct-tile premise \\
\lean{companion\_pose\_discrete} & A-C4.4 & \lean{companion\_pose\_discrete\_holds} & T1n & recovery of centre, normal, cubic frame, eighth-grid translation; census equality via list semantics; one side condition on fixed census metadata is the named T1n theorem \lean{mate\_records\_roles\_nonempty} (ruling R8) \\
\lean{only\_registered\_mates} & A-L5.3 & \lean{only\_registered\_mates\_holds} & T1 & endpoint keeps its frozen companion premises, including distinct placements; accepted companion/\allowbreak{}rigidity/\allowbreak{}census/\allowbreak{}overlap helpers assemble the literal 44-mate conclusion \\
\lean{component\_solids\_cover} & A-L6.2 & \lean{component\_solids\_cover\_holds} & T1 & clamp to an interior point with dist$^2$ $\le$ 3/\allowbreak{}2500 < 1/\allowbreak{}100, use the carried 1/\allowbreak{}4-ball and packing disjointness; alternative to assembly-wide tube gluing and conditional on an existing tiling, not an existence proof \\
\lean{small\_collar\_realization} & R\S{}8--9 & \lean{small\_collar\_realization\_holds} & T1n & duplicate tube sites are identified by centre; conjugated maps agree; membership equivalence holds for every tile including unrelated ones; separation is 21/\allowbreak{}200 > 1/\allowbreak{}16 \\
\lean{complete\_dihedral\_list} & A-L3.1 & \lean{complete\_dihedral\_list\_holds} & T1n & native wedge charts and all 6,408 mesh-edge incidences; the finite incidence trace is the named R8 theorem \lean{native\_mesh\_incidence\_trace} \\
\lean{generic\_feature\_partner} & A-L3.2 & \lean{generic\_feature\_partner\_holds} & T1n & native boundary strata and ordinary-meridian classification, with the finite exceptional-vertex tables certified by named native hooks \\
\lean{baseline\_component\_covers\_grid} & A-L6.1 & \lean{baseline\_component\_covers\_grid\_holds} & T1n & concrete companion assembly from the discharged local geometry and exact atlas \\
\lean{unrestricted\_alignment} & A-T6.3 & \lean{unrestricted\_alignment\_holds} & T1n & component coverage and the certified proper-frame group assemble global registration \\
\end{longtable}}


% generated by paper/scripts/axioms_table.py; from lean/R44/build_axioms.log at d90313a717; 169 declarations with an axiom line
{\small
\begin{longtable}{L{0.3\linewidth} L{0.58\linewidth}}
\caption{Axiom lines: the output of \lean{\#print axioms} for the finite theorems and the main theorems; ``standard'' is \lean{propext}, \lean{Classical.choice}, \lean{Quot.sound}, and hooks are the per-theorem \lean{native\_decide} axioms (FIGURES.md, Table 4).}\label{tab:axioms}\\
\toprule
declaration & axioms (\texttt{\#print axioms}) \\
\midrule
\endfirsthead
\toprule
declaration & axioms (\texttt{\#print axioms}) \\
\midrule
\endhead
\bottomrule
\endfoot
\lean{children\_partition\_2P} & none \\
\lean{contact\_closure\_30} & hooks: contact\_closure\_30 \\
\lean{profile\_canonical} & propext; hooks: profile\_canonical \\
\lean{atlas\_44} & standard; hooks: atlas\_44 \\
\lean{first\_shells\_33} & propext; hooks: first\_shells\_33 \\
\lean{central\_completion} & propext; hooks: central\_completion \\
\lean{parent\_conflicts\_28} & hooks: parent\_conflicts\_28 \\
\lean{parent\_atlas\_eq\_fine} & standard; hooks: parent\_atlas\_eq\_fine \\
\lean{mates\_census} & standard; hooks: mates\_census \\
\lean{deviations\_distinct} & none \\
\lean{no\_native\_symmetry} & propext \\
\lean{native\_features\_length} & propext \\
\lean{native\_feature\_centers\_nodup} & propext \\
\lean{native\_feature\_normal\_axis} & propext \\
\lean{solid\_mesh\_exact} & standard; hooks: solid\_mesh\_exact \\
\lean{boundary\_sphere} & standard; hooks: boundary\_sphere \\
\lean{nested\_substitution\_controls} & Quot.sound, propext; hooks: nested\_substitution\_controls \\
\lean{hierarchy\_cell\_controls} & hooks: hierarchy\_cell\_controls \\
\lean{registered\_shell\_cell\_controls} & propext; hooks: registered\_shell\_cell\_controls \\
\lean{orientation\_group\_24} & Quot.sound, propext; hooks: orientation\_group\_24 \\
\lean{substitution\_modular\_coincidence} & propext; hooks: substitution\_modular\_coincidence \\
\lean{seed\_language\_closure} & Quot.sound, propext; hooks: seed\_language\_closure \\
\lean{mesh\_angle\_audit} & standard; hooks: mesh\_angle\_audit \\
\lean{r44\_einstein} & standard; 21 hooks (DischargeCarrierStrata.carrier\_coordinate\_states\_table, DischargeMeshTrace.native\_mesh\_incidence\_trace, \dots) \\
\lean{tiling\_homochiral} & standard; 13 hooks (DischargeCarrierStrata.carrier\_coordinate\_states\_table, DischargeMeshTrace.native\_mesh\_incidence\_trace, \dots) \\
\lean{tiling\_chirality\_corollary} & standard; 13 hooks (DischargeCarrierStrata.carrier\_coordinate\_states\_table, DischargeMeshTrace.native\_mesh\_incidence\_trace, \dots) \\
\lean{carrier\_hierarchy} & standard; 17 hooks (DischargeCarrierStrata.carrier\_coordinate\_states\_table, DischargeMeshTrace.native\_mesh\_incidence\_trace, \dots) \\
\lean{carrier\_hierarchy\_unique} & standard; 17 hooks (DischargeCarrierStrata.carrier\_coordinate\_states\_table, DischargeMeshTrace.native\_mesh\_incidence\_trace, \dots) \\
\lean{substitution\_modular\_coincidence} & propext; hooks: substitution\_modular\_coincidence \\
\end{longtable}}


\subsection{The limit-periodic corollary}\label{sec:limitperiodic}

% ledger: L1, L2, L3, L4, L5, L6, L7, Q15, Q16, Q17, Q18, Q19, N1; DECISIONS Q-A (fixed-point wording), Q-B (seed); Lean substitution_modular_coincidence (T1n); verify/substitution_modular_coincidence.py (T2); paper/scripts/seed_language_closure.py (T2)
The substitution of \cref{sec:substitution} is a lattice substitution in the sense of Lee and
Moody \cite{LM01}. Label each unit cell of a registered tiling by the frame of the chair that
covers it and by which of the seven cells of that chair it is: $24\times7=168$ labels. The
refinement replaces each labelled cell by its eight labelled children at twice the scale, so the
labelling is a substitution on $\Z^3$ with inflation $2$ and $168$ symbols, and its $168\times168$
substitution matrix has constant column sum $8=|\det 2I|$.

\begin{theorem}[modular coincidence; \tier{T1n}, replayed in Python]\label{thm:coincidence}
The $168$-label substitution is total and primitive, with least primitive exponent $N=3$, and it
has a modular coincidence of least depth $M=3$: at the address $a=(0,0,2)$ modulo $2^3$, the
level-$3$ descendant of every one of the $168$ labels has label $i=78$.
\end{theorem}

The Lean theorem is \lean{substitution\_modular\_coincidence}. The check tracks the full set of
$168$ labels from the root, so the coincidence is over all labels, not along a single orbit; the integers $N$ and $M$ are independent and happen to
coincide.

% ledger: L3, L4, L9, L10, N1, T9; DECISIONS Q-B; Lean seed_language_closure (T1n); verify/replay.py step seed_language_closure (T2); paper/scripts/seed_language_closure.py (T2); external review of the manuscript 2026-09-09, findings B1 and M1 (paper/review/EXTERNAL_RULING_2026-09-09.md)
A legal fixed point on the whole of $\Z^3$ exists. A substitution-legal fixed labelling of
$\Z^3$ is determined by its seed on $\{-1,0\}^3$: the cell $-\delta$, $\delta\in\{0,1\}^3$, has
parent $-\delta$ and residue $\delta$, so its eight labels satisfy their fixed-residue equations;
conversely, a seed in the substitution language satisfying those equations yields such a
labelling, whereas a block satisfying the equations but outside the language need not (one such
block encodes inconsistent whole-chair labels). The substitution language
contains $27$ fixed seed blocks on $\{-1,0\}^3$, in proper-rotation orbits of sizes $24$ and $3$;
iterating any one gives a legal substitution-fixed labelling of $\Z^3$ (\cref{fig:coincidence}).
The count is by exact closure of the language of legal $2\times2\times2$ blocks, which grows
through $168$, $600$, $1{,}278$, $1{,}398$ and $1{,}410$ blocks under successive substitution and
then is stable, and every seed is checked to reproduce itself in place under iteration (the
theorem \lean{seed\_language\_closure}, \tier{T1n}, whose statement lists the five closure sizes,
the count $27$, the orbit sizes $24$ and $3$, the stabilizer histogram, in-place reproduction and
commutation with the $24$ frames; the replay step of the same name in \file{verify/replay.py};
and, independently, \file{paper/scripts/seed\_language\_closure.py}, standard library,
\cref{app:data}); a bounded
search inside level-$3$ patches, which an earlier draft reported as the count $24$, misses the
three seeds of the small orbit, and the external review of the manuscript found them. Every
finite patch of a fixed labelling occurs in some $\sigma^n(j)$ and hence, by primitivity, inside
a refinement of the native chair, that is, inside the tiling of \cref{thm:existence}; so the
chair poses the labelling encodes form a registered chair tiling of $\R^3$ all of whose contacts
lie in $\atlas$, and small-collar realization (\cref{sec:existence}) turns it into a tiling by
$\solid$ at those poses. We write $T_w$ for the tiling determined by the seed $w$.

\begin{proposition}[symmetries of the fixed points; finite data \tier{T1n}
(\lean{seed\_language\_closure}) and \tier{T2}, the deduction
written]\label[proposition]{prop:seedsym}
Let $G$ be the $24$-element proper cubic group, acting on the seed cube by moving the cells and
composing the frames in the labels. For each of the $27$ fixed seeds $w$, $\Sym(T_w)$ is the
stabilizer of $w$ in $G$. The stabilizers have order $1$ for the $24$ seeds of the large orbit
and order $8$ for the three seeds of the small orbit; for the seed $(28,84,91,35,98,42,49,105)$,
in the repository's label order, the stabilizer is generated by $(x,y,z)\mapsto(z,y,-x)$ and
$(x,y,z)\mapsto(-x,-y,z)$. Hence some tilings by $\solid$ have trivial symmetry group and some
have symmetry group of order $8$.
\end{proposition}

\begin{proof}
Let $g=(R,t)$, $x\mapsto Rx+t$, be a symmetry of $T_w$. As in the proof of \cref{thm:sym},
registration and the absence of self-isometries give $R\in G$ and $t\in\Z^3$. The parent
partition is read from the tiling alone, and the children of a rotated pose are the rotated
children (a finite check over $G$ in the same script), so by the uniqueness of parents
(\cref{thm:parent}) the parents of $RT$ are the rotated parents of $T$; with the translation
clause of \cref{thm:coarsening}, $D(gT)=(R,t/2)\,D(T)$. The seed being fixed in place,
$D(T_w)=T_w$, so $(R,t/2)$ is again a symmetry of $T_w$ and $t/2\in\Z^3$; iterating,
$t\in\bigcap_n2^n\Z^3=\{0\}$. Thus $\Sym(T_w)=\{R\in G: RT_w=T_w\}$. The substitution commutes
with the action of $G$ on labelled cells (the same finite check), so $RT_w$ is the fixed point
with seed $Rw$; a fixed point is determined by its seed, so $RT_w=T_w$ exactly when $Rw=w$. The
orders and the generators are read off the $27$ seeds.
\end{proof}

The seed $(6,0,4,5,2,3,0,1)$ of the large orbit is the one whose ancestors fill only an octant
(\cref{sec:hierarchy}); its tiling has trivial symmetry group. The three seeds of the small orbit
each consist of eight whole tiles at the origin, one in each octant, each its own child $000$ at
every level.

% ledger: D9, Q16; source math0002019 L102-L107 (Lee-Moody definitions)
We recall the definitions the corollary uses, from Lee and Moody \cite{LM01}. A cut-and-project
scheme consists of a lattice $\widetilde L\subset\R^3\times G$, where $G$ is a locally compact
abelian group, such that the physical projection $\pi_1$ is injective on $\widetilde L$ and the
internal projection $\pi_2$ has dense image in $G$. A model set is, up to translation,
$\{\pi_1(\ell):\ell\in\widetilde L,\ \pi_2(\ell)\in W\}$, where the window $W\subset G$ has
nonempty interior and compact closure; it is regular when $\partial W$ has Haar measure zero.
Here the internal group is $\varprojlim_n\Z^3/2^n\Z^3$, the $2$-adic completion. A modular
coincidence occurs when, at some substitution depth and residue address, all substitution maps
at that address have the same output label.

% ledger: L1, L2, L3, L4, L5, L6, L7, Q15, Q16, Q17, Q18, Q19, N1; DECISIONS Q-A (fixed-point wording); the corollary and the paragraph after it
\begin{corollary}[limit-periodic structure of the fixed point; the cited imports are \tier{T3}]\label[corollary]{cor:limitperiodic}
By Lee and Moody's Theorem 3 \cite{LM01}, applied in the direction (iv)~$\Rightarrow$~(ii) with
its hypotheses verified as above (\cref{thm:coincidence}: primitivity; Perron--Frobenius eigenvalue $8=|\det 2I|$; the
label classes of the fixed point partition $\Z^3$), the $168$ label classes of the substitution
fixed point are regular model sets for the $2$-adic cut-and-project scheme of \cite{LM01}; by
Schlottmann's theorem, in the form stated by Lee, Moody and Solomyak
\cite[Theorem~5.11]{LMS03}, each is pure point diffractive; in the terminology of Baake, Moody and
Schlottmann and of Baake and Grimm \cite{BMS98,BG10,BG13} the structure is limit-periodic.
\end{corollary}

The corollary is a statement about the fixed point, not about every tiling by $\solid$: whether
every tiling lies in the substitution hull is the fault-line question left open in
\cref{sec:hierarchy}, and Lee and Moody's own passage from the substitution to every tiling
obeying the Taylor--Socolar matching rules is a separate three-part argument \cite{LM13}, not
an application of their Theorem 3. In Socolar's words about the hat, whose structure ``is
qualitatively different from the limit-periodic structure of the hexagonal Taylor--Socolar
monotile tiling, which has Bragg diffraction peaks at wavelengths of the form $2^{-n}k_0$ for
arbitrarily large $n$'' \cite{Soc23}, the fixed-point label-class structure established here is
limit-periodic with a $2$-adic internal space, rather than quasicrystalline in the sense
contrasted by Socolar.
Akiyama and Lee verified pure point spectrum for a $168$-prototile Taylor--Socolar substitution by
an overlap coincidence \cite{AL14}, and Lee and Moody checked the sphinx's modular coincidence by
computer at depth $8$ \cite{LM01}; ours is checked at depth $3$ in Lean and replayed in Python.

\begin{figure}[ht]
  \centering
  % generated by paper/scripts/coincidence_tree.py from verify/substitution_modular_coincidence.py
\begin{tikzpicture}[x=1cm,y=1cm,font=\small,
  lvl/.style={draw,rounded corners=1pt,inner sep=3pt,align=center,minimum width=2.6cm},
  arr/.style={-{Latex[length=1.6mm]},thick}]
  \node[lvl] (d0) at (0,0) {root: all $168$ labels};
  \node[lvl] (d1) at (0,-1.5) {residue $(0,0,0)$ at depth 1\\ \scriptsize all-address histogram: 8\,addresses with 42};
  \draw[arr] (d0) -- (d1);
  \node[lvl] (d2) at (0,-3.0) {residue $(0,0,1)$ at depth 2\\ \scriptsize all-address histogram: 48\,addresses with 6, 16\,addresses with 24};
  \draw[arr] (d1) -- (d2);
  \node[lvl] (d3) at (0,-4.5) {residue $(0,0,0)$ at depth 3\\ \scriptsize all-address histogram: 336\,addresses with 1, 160\,addresses with 6, 16\,addresses with 24};
  \draw[arr] (d2) -- (d3);
  \node[lvl,thick,fill=black!6] (fin) at (0,-6.0) {address $a=(0,0,2)$ mod $2^{3}$: one label, $i=78$};
  \draw[arr] (d3) -- (fin);
\end{tikzpicture}

  \caption{The modular-coincidence address and the all-address distributions of label-set sizes
  at depths $1$--$3$.}
  \label{fig:coincidence}
\end{figure}
